% GENERATED FILE. Edit canonical lessons, then run npm run generate:books.
% canonical-source-hash: fnv1a64:688565bbe0abdc16
% canonical-lessons: ZH-C523-pifa, ZH-C523-maimai, ZH-W523-mai-sell, ZH-C523-maijia-seller, ZH-C523-maijia-buyer, ZH-C523-lingshou

\chapter{Buying and Selling}
\label{ch:zh-buying-and-selling}
\hlchaptermodality{523}

\begin{chapteropening}
\textbf{By the end of this chapter:} \emph{I can talk about running a business, sellers and buyers, and selling wholesale or retail.}

Complete the last lesson of chapter 523: I can talk about running a business, sellers and buyers, and selling wholesale or retail.
\end{chapteropening}

\section[pīfā]{pīfā — to sell wholesale}

\label{lesson:ZH-C523-pifa}



\begin{quote}
Before the new one: say the Mandarin for a doorbell, then the Mandarin for the floor of a room.
\end{quote}

\subsection*{You'll want to know: pīfā}
\emph{pīfā} — \textquotedblleft{}to sell wholesale\textquotedblright{}.

\begin{quote}
\emph{zhè jiā shāngdiàn pīfā shuǐguǒ} — This shop sells fruit wholesale.
\end{quote}

Its characters come later; for now it is a word for the ear.

\begin{cousinweb}
\emph{pīfā} is built from \emph{pī} and \zh{发} 'to send out', both inherited from Middle Chinese, \emph{phej} and \emph{pjot} (Guangyun, in Baxter's transcription). The dictionary glosses the first 'to strike', 'to push' and 'to pass on'; how the pair came to mean selling in bulk is not settled.
\end{cousinweb}

\begin{grammarlens}[title={What you've built}]
One more word for saying how it was.
\end{grammarlens}

\subsection*{Your turn}
\begin{itemize}
\raggedright
  \item \emph{Say it:} \emph{pīfā}
  \item \emph{Say it:} \emph{pīfā}, once more
  \item \emph{Say it:} \emph{zhè jiā shāngdiàn pīfā shuǐguǒ} — This shop sells fruit wholesale.
  \item \emph{Recall:} say the Mandarin for a doorbell, then the Mandarin for the floor of a room, then say \emph{pīfā} again
\end{itemize}

\begin{tcolorbox}[breakable,colback=teal!4,colframe=teal!35!black,title={Before you move on}]
What does pīfā mean? (\textquotedblleft{}To sell wholesale\textquotedblright{}.) Say it once more.
\end{tcolorbox}
\section[\texorpdfstring{\zh{买卖} (mǎimai)}{买卖 (mǎimai)}]{\zh{买卖} (mǎimai) — buying and selling for a living; a business deal}

\label{lesson:ZH-C523-maimai}



\begin{quote}
Before the new one: say the Mandarin for the floor of a room, then the Mandarin for to sell wholesale.
\end{quote}

\subsection*{You'll want to know: \zh{买卖}}
\textbf{\zh{买卖}} — \emph{mǎimai} — \textquotedblleft{}buying and selling for a living; a business deal\textquotedblright{}.

\begin{quote}
\textbf{\zh{他做了二十年买卖}} \emph{tā zuò le èrshí nián mǎimai} — He has made his living buying and selling for twenty years.
\end{quote}

The new character, in pieces: \zh{卖} — eight strokes in two pieces: \zh{十} on top, two strokes; and \zh{买} under it, six strokes. The next lesson puts it in your hand.

\begin{cousinweb}
\textbf{\zh{买卖}} (\emph{mǎimai}) is built from \zh{买} 'to buy' and \zh{卖} 'to sell', both inherited from Middle Chinese, \emph{meaX} and \emph{meaH} (Guangyun, in Baxter's transcription), which glosses the first, after an older dictionary, 'to trade' and the second 'to put goods out'. The two differ only in their tones, and in the word the second is said light.
\end{cousinweb}

\begin{grammarlens}[title={What you've built}]
One more word for saying how it was.
\end{grammarlens}

\subsection*{Your turn}
\begin{itemize}
\raggedright
  \item \emph{Say it:} \emph{mǎimai}
  \item \emph{Say it:} \emph{mǎimai}, once more
  \item \emph{Say it:} \emph{tā zuò le èrshí nián mǎimai} — He has made his living buying and selling for twenty years.
  \item \emph{Recall:} say the Mandarin for the floor of a room, then the Mandarin for to sell wholesale, then say \emph{mǎimai} again
\end{itemize}

\begin{tcolorbox}[breakable,colback=teal!4,colframe=teal!35!black,title={Before you move on}]
What does \zh{买卖} mean? (\textquotedblleft{}Buying and selling for a living; a business deal\textquotedblright{}.) Say it once more.
\end{tcolorbox}
\section[\texorpdfstring{\zh{卖} (mài)}{卖 (mài)}]{\zh{卖} — \zh{十} on top, then \zh{买} under it}

\label{lesson:ZH-W523-mai-sell}



\begin{quote}
The word first, then the character inside it.

\begin{itemize}
\raggedright
  \item \emph{Recall:} say \emph{mǎimai} — \textbf{R1}, one lesson back, before it decays
\end{itemize}

\begin{quote}
\zh{买} \zh{十}
\end{quote}

\emph{Read it:} each one aloud, with nothing said first — \emph{mǎi}, \emph{shí}
\end{quote}

\subsection*{Script — \zh{十} on top, then \zh{买} under it}
\hlblockfigure{\detokenize{figures/ZH-W523-mai-sell-filmstrip.pdf}}{How \zh{卖} is written, stroke by stroke}

\begin{quote}
\zh{卖}
\end{quote}

\textbf{Two pieces, top to bottom.} Count the pieces, not the strokes.

\textbf{The top piece, \zh{十}, two strokes}: a short bar, then a short vertical down through it.

\textbf{The bottom piece, \zh{买}, six strokes}, as you wrote it: a bar that hooks down at its right end, two short strokes falling to the right under it, a long bar, a long stroke falling to the left and a long stroke falling to the right.

2 + 6 = eight: to sell is to buy with \zh{十} on top.

\subsection*{Your turn}
\begin{itemize}
\raggedright
  \item \emph{Write it:} \zh{卖}, all eight, with the model in front of you
  \item \emph{Cover:} write it again from memory, and count the pieces, not the strokes
\end{itemize}

\begin{tcolorbox}[breakable,colback=teal!4,colframe=teal!35!black,title={Before you move on}]
How many pieces does \zh{卖} have? (\textbf{Two: \zh{十} and \zh{买}}.) How does \zh{卖} differ from \zh{买}? (\textbf{It has \zh{十} on top}.)

Source: \href{https://raw.githubusercontent.com/chanind/hanzi-writer-data/68d10a4b21150cae5e1ebbd223eed289cf32d90c/data/\%E5\%8D\%96.json}{Hanzi Writer Data \zh{卖}, pinned stroke paths and medians}.
\end{tcolorbox}
\section[\texorpdfstring{\zh{卖家} (màijiā)}{卖家 (màijiā)}]{\zh{卖家} (màijiā) — seller}

\label{lesson:ZH-C523-maijia-seller}



\begin{quote}
Before the new one: say the Mandarin for to sell wholesale, then the Mandarin for buying and selling for a living.
\end{quote}

\subsection*{You'll want to know: \zh{卖家}}
\textbf{\zh{卖家}} — \emph{màijiā} — \textquotedblleft{}seller\textquotedblright{}.

\begin{quote}
\textbf{\zh{我觉得这个卖家最好}} \emph{wǒ juéde zhège màijiā zuì hǎo} — I think this seller is the best.
\end{quote}

\begin{cousinweb}
\textbf{\zh{卖家}} (\emph{màijiā}) is built from \zh{卖} 'to sell' and \zh{家}, both inherited from Middle Chinese, \emph{meaH} and \emph{kae} (Guangyun, in Baxter's transcription). The dictionary glosses the second 'a dwelling'; its use for the one who does a trade is not given there.
\end{cousinweb}

\begin{grammarlens}[title={What you've built}]
One more word for saying how it was.
\end{grammarlens}

\subsection*{Your turn}
\begin{itemize}
\raggedright
  \item \emph{Say it:} \emph{màijiā}
  \item \emph{Say it:} \emph{màijiā}, once more
  \item \emph{Say it:} \emph{wǒ juéde zhège màijiā zuì hǎo} — I think this seller is the best.
  \item \emph{Recall:} say the Mandarin for to sell wholesale, then the Mandarin for buying and selling for a living, then say \emph{màijiā} again
\end{itemize}

\begin{tcolorbox}[breakable,colback=teal!4,colframe=teal!35!black,title={Before you move on}]
What does \zh{卖家} mean? (\textquotedblleft{}Seller\textquotedblright{}.) Say it once more.
\end{tcolorbox}
\section[\texorpdfstring{\zh{买家} (mǎijiā)}{买家 (mǎijiā)}]{\zh{买家} (mǎijiā) — buyer}

\label{lesson:ZH-C523-maijia-buyer}



\begin{quote}
Before the new one: say the Mandarin for buying and selling for a living, then the Mandarin for the seller (in a sale).
\end{quote}

\subsection*{You'll want to know: \zh{买家}}
\textbf{\zh{买家}} — \emph{mǎijiā} — \textquotedblleft{}buyer\textquotedblright{}.

\begin{quote}
\textbf{\zh{买家都觉得这个卖家很好}} \emph{mǎijiā dōu juéde zhège màijiā hěn hǎo} — The buyers all think this seller is very good.
\end{quote}

\begin{cousinweb}
\textbf{\zh{买家}} (\emph{mǎijiā}) is built from \zh{买} 'to buy' and \zh{家}, both inherited from Middle Chinese, \emph{meaX} and \emph{kae} (Guangyun, in Baxter's transcription), which glosses the second 'a dwelling'.
\end{cousinweb}

\begin{grammarlens}[title={What you've built}]
One more word for saying how it was.
\end{grammarlens}

\subsection*{Your turn}
\begin{itemize}
\raggedright
  \item \emph{Say it:} \emph{mǎijiā}
  \item \emph{Say it:} \emph{mǎijiā}, once more
  \item \emph{Say it:} \emph{mǎijiā dōu juéde zhège màijiā hěn hǎo} — The buyers all think this seller is very good.
  \item \emph{Recall:} say the Mandarin for buying and selling for a living, then the Mandarin for the seller (in a sale), then say \emph{mǎijiā} again
\end{itemize}

\begin{tcolorbox}[breakable,colback=teal!4,colframe=teal!35!black,title={Before you move on}]
What does \zh{买家} mean? (\textquotedblleft{}Buyer\textquotedblright{}.) Say it once more.
\end{tcolorbox}
\section[língshòu]{língshòu — to sell retail}

\label{lesson:ZH-C523-lingshou}



\begin{quote}
Before the new one: say the Mandarin for the seller (in a sale), then the Mandarin for the buyer (in a sale).
\end{quote}

\subsection*{You'll want to know: língshòu}
\emph{língshòu} — \textquotedblleft{}to sell retail\textquotedblright{}.

\begin{quote}
\emph{zhè jiā shāngdiàn zhǐ língshòu bù pīfā} — This shop sells only retail, not wholesale.
\end{quote}

Its characters come later; for now it is a word for the ear.

\begin{cousinweb}
\emph{língshòu} is built from \emph{líng} and \emph{shòu} 'to sell', both inherited from Middle Chinese, \emph{leng} and \emph{dzyuwH} (Guangyun, in Baxter's transcription). The dictionary glosses the first 'to fall' and, after an older dictionary, as a slow rain; the sense 'in small bits' is not given there. It glosses the second 'to sell goods'.
\end{cousinweb}

\begin{grammarlens}[title={What you've built}]
One more word for saying how it was.
\end{grammarlens}

\subsection*{Your turn}
\begin{itemize}
\raggedright
  \item \emph{Say it:} \emph{língshòu}
  \item \emph{Say it:} \emph{língshòu}, once more
  \item \emph{Say it:} \emph{zhè jiā shāngdiàn zhǐ língshòu bù pīfā} — This shop sells only retail, not wholesale.
  \item \emph{Recall:} say the Mandarin for the seller (in a sale), then the Mandarin for the buyer (in a sale), then say \emph{língshòu} again
\end{itemize}

\begin{tcolorbox}[breakable,colback=teal!4,colframe=teal!35!black,title={Before you move on}]
What does língshòu mean? (\textquotedblleft{}To sell retail\textquotedblright{}.) Say it once more.
\end{tcolorbox}
